\begin{figure}
  \centering
  \includegraphics[width=\linewidth]{tornado.pdf}
  \caption{Cost change needed to move the optimal schedule (100~MW new solar, minimising production cost). Each bar spans the multipliers on that technology's production cost that leave the optimal schedule unchanged; its ends are the breakevens, annotated as percentage changes. The axis is logarithmic, so that halving and doubling a cost are equally far from the centre. A bar running to the plot edge, marked ``no flip'', means the schedule does not change in that direction out to a factor of 10: a technology that is already fully dispatched is unaffected by becoming cheaper, and one that is switched off is unaffected by becoming dearer. Technologies are ordered by their nearer breakeven, most fragile first. The breakevens are identical in all four solar scenarios.}
  \label{fig:tornado}
\end{figure}

\begin{figure}
  \centering
  \includegraphics[width=\linewidth]{flip_consequences.pdf}
  \caption{What the first flip is worth. For each technology and each direction, the change in total emissions once the schedule has flipped at the breakeven of Figure~\ref{fig:tornado}, evaluated at true cost coefficients, against the amount of new solar capacity built. Left: the technology becomes cheaper. Right: it becomes dearer. Absent series have no breakeven in that direction. Read together with Figure~\ref{fig:tornado}: a breakeven that is close by is not necessarily one that matters.}
  \label{fig:consequences}
\end{figure}
